\section{Repaso del curso y contexto del problema}

Esta es la clase 12 de la serie de KAIST CS492D sobre ``modelos de difusión'', que continúa con el contenido de la sesión anterior sobre generación guiada en tiempo de inferencia (Inference-Time Guided Generation), profundizando en técnicas de guía más avanzadas. La clase anterior presentó tres categorías principales de métodos de guía en tiempo de inferencia:

\begin{enumerate}
    \item \textbf{Manipulación de atención y capas intermedias}: lograr la guía mediante la modificación directa de las capas intermedias o las salidas de atención de la red neuronal. La ventaja es que es fácil de implementar, pero la desventaja es que carece de una explicación teórica y solo funciona en escenarios específicos.
    \item \textbf{Guía basada en problemas inversos (DPS)}: modelar la generación condicional como un problema inverso, descomponiendo el puntaje condicional usando la regla de Bayes en la suma del puntaje marginal y el gradiente de verosimilitud condicional.
    \item \textbf{Guía basada en funciones de recompensa}: definir las preferencias del usuario mediante una función de recompensa diferenciable, inyectando guías de gradiente durante el proceso de generación.
\end{enumerate}

\begin{importantbox}{Repaso de problemas inversos}
Dado un observable $y$ y una función de mapeo $\mathcal{A}$, el objetivo del problema inverso es encontrar la $x$ que satisfaga $y = \mathcal{A}(x)$. En el marco de los modelos de difusión, necesitamos calcular el puntaje condicional $\nabla_{x_t} \log p(x_t | y)$, descomponiéndolo usando la regla de Bayes en:
$$\nabla_{x_t} \log p(x_t | y) = \nabla_{x_t} \log p(x_t) + \nabla_{x_t} \log p(y | x_t)$$
donde el primer término es el puntaje marginal que ya aprendió el modelo preentrenado, y el segundo término es el gradiente de verosimilitud condicional que requiere un cálculo adicional.
\end{importantbox}

\subsection{Aproximación central del método DPS}

La idea clave del método DPS (Diffusion Posterior Sampling) es: cuando la función de mapeo $\mathcal{A}$ en el problema inverso es lineal, el término puntaje condicional puede calcularse mediante el parámetro $R$ (la pseudoinversa de $\mathcal{A}$). Además, se puede obtener una forma más concisa mediante una aproximación más gruesa: aproximar la distribución condicional como una distribución gaussiana:

$$\nabla_{x_t} \log p(y | x_t) \approx -\nabla_{x_t} \| y - \mathcal{A}(\hat{x}_0(x_t)) \|^2$$

donde $\hat{x}_0(x_t)$ es los datos limpios estimados a partir del estado ruidoso actual $x_t$ usando la fórmula de Tweedie.

\begin{knowledgebox}{Fórmula de Tweedie}
La fórmula de Tweedie proporciona la estimación posterior óptima de los datos limpios $x_0$ dados una observación ruidosa $x_t$:
$$\hat{x}_0(x_t) = \frac{1}{\sqrt{\bar{\alpha}_t}} \left( x_t + (1 - \bar{\alpha}_t) \nabla_{x_t} \log p(x_t) \right)$$
En la práctica, esto equivale a estimar $x_0$ usando una red de predicción de ruido preentrenada $\epsilon_\theta(x_t, t)$.
\end{knowledgebox}

\subsection{Simplicidad en la implementación de DPS}

El método DPS es extremadamente sencillo de implementar: solo requiere agregar una línea de código al pipeline de generación. En cada paso de desruido, primero se calcula $x_{t-1}$ mediante el proceso inverso estándar y luego se realiza una actualización de un paso a $x_{t-1}$ basándose en el gradiente de la pérdida L2:

$$x_{t-1} \leftarrow x_{t-1} - \eta \cdot \nabla_{x_t} \| y - \mathcal{A}(\hat{x}_0(x_t)) \|^2$$

donde $\eta$ es el tamaño del paso de guía (guidance step size), que controla la intensidad de la guía en cada paso. Todo el flujo del algoritmo DPS puede resumirse así:

\begin{enumerate}
    \item Comenzar con $x_T \sim \mathcal{N}(0, I)$
    \item Para cada paso temporal $t = T, T-1, \ldots, 1$:
    \begin{enumerate}
        \item Calcular $\hat{x}_0(x_t)$ usando el modelo preentrenado (estimación de Tweedie)
        \item Calcular el gradiente de guía $g = \nabla_{x_t} \| y - \mathcal{A}(\hat{x}_0(x_t)) \|^2$
        \item Ejecutar el paso de desruido estándar para obtener $\tilde{x}_{t-1}$
        \item Aplicar la actualización de guía $x_{t-1} = \tilde{x}_{t-1} - \eta_t \cdot g$
    \end{enumerate}
    \item Salir con $x_0$
\end{enumerate}

El ponente recomienda encarecidamente a los estudiantes que experimenten el método DPS usando su propio código de modelos de difusión (como la implementación en las tareas del curso), ya que realmente solo requiere agregar una línea de código de actualización de gradiente al pipeline de generación.

\begin{warningbox}{Costo de la aproximación}
Aunque aproximar $p(y|x_t)$ como una distribución gaussiana hace que los cálculos sean extremadamente simples, esto no representa la distribución real. Especialmente en pasos temporales con alto nivel de ruido, la calidad de la estimación de $\hat{x}_0$ es deficiente y el error de aproximación de la verosimilitud condicional puede ser muy grande. Por lo tanto, el rendimiento práctico suele ser muy sensible al peso de guía y al número de pasos temporales.
\end{warningbox}

\begin{knowledgebox}{Problemas inversos lineales vs. no lineales}
Cuando la función de mapeo $\mathcal{A}$ es lineal (como en superresolución, desenfocamiento o reparación), se puede derivar una expresión más precisa del puntaje condicional que involucra operaciones de pseudoinversa de $\mathcal{A}$. Sin embargo, en la práctica, incluso en casos lineales, la aproximación gaussiana simple de DPS logra buenos resultados; por lo tanto, generalmente se usa directamente la forma general de guía por gradiente L2. Para mapeos no lineales $\mathcal{A}$ (como recuperación de fase o modelos de sensores no lineales), el error de la aproximación gaussiana puede ser mayor, pero sigue siendo el método más práctico hasta ahora.
\end{knowledgebox}

\subsection{Resumen de la sección}

El marco DPS descompone el puntaje condicional en un puntaje marginal y un gradiente de verosimilitud condicional mediante la regla de Bayes, y simplifica este último usando una aproximación gaussiana a un gradiente de pérdida L2. Este método no requiere entrenamiento adicional y puede aplicarse a cualquier función condicional diferenciable.
